% Design sample: one chapter opener page with the core elements.
% Compile from the project root:  bash scripts/build.sh style/sample.tex  ->  build/sample.pdf
% Placeholder content only; numbers here are not report data.
\documentclass{style/report}

\begin{document}
\setcounter{chapter}{2}
\chapter{סקירת ספרות}

\section{עקביות חזותית של דמויות}

פסקה זו מדגימה את גוף הטקסט: גופן \textenglish{David} בגודל 12 נקודות, ריווח שורות 1.5 ויישור דו־צדדי.
מערכות ליצירת וידאו (\textenglish{Text-to-Video}) מתמקדות בעקביות חזותית של הדמות בין סצנות,
ושיטות כגון \textenglish{IP-Adapter} ו־\textenglish{LoRA} משמשות לשמירה על מראה הפנים.
המספרים בטבלה \ref{tab:sample} הם דוגמה לעיצוב בלבד [TO-VERIFY].

\begin{table}[h]
  \centering
  \caption{דוגמה לטבלה בעיצוב הדוח}
  \label{tab:sample}
  \begin{tabular}{l l c}
    \headrow \hcell{כלי} & \hcell{סוג} & \hcell{תמיכה בעברית} \\
    \textenglish{Tool A} & מסחרי     & חלקית \\ \rowrule
    \textenglish{Tool B} & קוד פתוח  & לא    \\ \rowrule
    \textenglish{Tool C} & מסחרי     & מלאה  \\
  \end{tabular}
\end{table}

\begin{directornote}
  בפרויקט שלנו, עקביות הדמות נמדדת על פני לפחות עשר סצנות, ובנוסף נבדקת רציפות העלילה בין פרקים.
\end{directornote}

\begin{notebox}
  תיבת הערה כללית: רקע בצבע \textenglish{brandHighlight} וטקסט בצבע הדיו.
\end{notebox}

\end{document}
